\paragraph{IV 20.10: pencils of hyperplane sections and blown-up varieties (cote 25, sites 58--62, opening).}
Paragraph 10.1 relates the constructible exceptional set for a property of hyperplane sections to lines in the dual projective space, and defines the pencil associated with such a line. The discussion uses the results of sections 8 and 9, cites 8.2 for the additional regularity condition, and refers to 7.1 and paragraph 5 in its passage from a field to a neighborhood on a general base. The handwritten reference at site 58 reads 12 over the canceled word ``suivant''. Paragraph 10.2 begins at site 59 with the polar linear subvariety $L^o$ of codimension two and the fiber product $T=X\mathbin{\mathrm x}_P L^o$. It describes $T$ through the homomorphism from the rank-two submodule and through its ideal of zeros. The regularity hypothesis is expressed using a totally ordered local basis, then an isomorphism with $J/J^2$ together with a regular immersion. Th\'eor\`eme 10.2 states the unique $X$-isomorphism between the pencil and the blow-up of $X$ along $T$. The handwritten uniqueness phrase and the tensor insertion in the preceding condition are retained with their actual attachment points.

Site 60 discusses how to organize that theorem and the proposed statement 9.5, asks for the regularity requirement in the recapitulative statements, and describes the intended fibrations over $\mathbf P^1$. Its corrected paragraph reference is again 12; the attached note has an unresolved first character in S.3./5.3. The words ``sont lisses sont au plus'' are printed in that order and remain in the edition. Site 61 gives the more general Proposition 10.3, formulated for a quasi-coherent Module, its homomorphism to $\underline O_X$, and the induced construction over $P(\underline G)$. The following note calls its first assertion nearly immediate and leaves the second as an exercise for Dieudonn\'e. Two large crossing diagonals traverse this proposition and note, while the preceding paragraph suggests moving the material to the discussion of regular immersions. Their editorial function remains unresolved, so the entire block is retained with the graphic marks recorded. Four number groups on this page clearly show handwritten 10 over typed 8; the ambiguous connector in its displayed sequence is preserved without completing it. The final paragraph combines 10.2 with 9.2 in dimension at most one and leads to Corollaire 10.4 for a smooth geometrically connected curve in projective space over an infinite field. The corollary continues at the opening of site 62: the degree is revised by hand to $d$, and the fiber description contains $d-2$ copies of the indicated spectrum together with $\operatorname{Spec}k'[t]/(t^2)$. Its following explanation uses the term ``quadratique'' for the ramification. The old canceled fiber description and its uncertain coefficient remain in the apparatus.
